\documentclass[10pt,a4paper]{amsart}
\usepackage[margin=19mm]{geometry}
\usepackage{amsmath,amssymb,mathtools}
\usepackage[scr=rsfs,cal=euler]{mathalfa}
\usepackage{xfrac}
\usepackage[unicode,hidelinks,bookmarks=false]{hyperref}
\newcommand{\ie}{i.e.\ }
\newcommand{\cf}{cf.\ }
\newcommand{\resp}{respectively\ }
\newcommand{\Subsection}{\subsection}
\providecommand{\nobreakdash}{\nobreak\hbox{-}}
\setlength{\emergencystretch}{2em}
\multlinegap0pt
\usepackage{amssymb}
\usepackage{tikz}
\usepackage{mathtools}
\usepackage{xcolor}
\usepackage{enumerate}
\newtheorem{lemma}{Lemma}
\title{Dice Relabeling using Square-Sided Dice}
\author{Evelyn Fiore, George D. Nasr, Cooper Stone}
\date{Source: arXiv:2606.20311v1 (CC BY 4.0)}
\begin{document}
\maketitle

\noindent\emph{Source and licence.} Dice Relabeling using Square-Sided Dice. Version arXiv:2606.20311v1 (CC BY 4.0), distributed by arXiv under the Creative Commons Attribution 4.0 International licence, http://creativecommons.org/licenses/by/4.0/, as recorded in the arXiv licence metadata for this exact version and shown on the abstract page https://arxiv.org/abs/2606.20311v1. Full licence notice: \texttt{licenses/arxiv-2606.20311-CC-BY-4.0.txt}.\par\smallskip

\noindent\textit{Editorial scope.} Counted unit: the article's lemma lem:q01\_p2\_comp (source0.tex lines 1317--1325). The article's complete original proof of it is delivered with it (source0.tex lines 1326--1396, one \texttt{proof} environment of 71 lines). No mathematics is written, repaired or re-derived: the statement and the proof are the article's own lines, in the article's own notation, and no retained line is substituted or re-flowed.  No further source-local context is needed: every cross-reference of every retained line resolves inside the two delivered spans.  Nothing was omitted from any retained span, so the package carries no omission record.  No retained line contains a citation, so the wrapper carries no bibliography and no bibliography entry.  No retained line contains a figure environment, an image-inclusion command or drawing code, so no image is delivered and the package declares no asset dependency.  Editorial formatting only, in addition: an AMS \texttt{amsart} preamble that restates the article's own declarations and macros used by the retained lines, namely the environments \texttt{lemma} on separate counters, together with the standard packages those lines need.  The retained environments print in this document as Lemma 1 in order of appearance, whereas the article prints the same items with the numbering of its own full document.  The complete native source of the article as distributed in the arXiv e-print for this exact version is bundled unchanged as \texttt{sources/203206/source0.tex}.
\begin{lemma}\label{lem:q01_p2_comp}
Let $\alpha=\lfloor {k\over 2}\rfloor$ and $\beta=\lfloor{k-q\over 2}\rfloor$. Then 

\begin{equation}\label{eq:floorsnonneg}
\delta(\alpha \text{ even}) k -(-1)^{\alpha}\alpha +\delta(\beta \text{ even})(2q-2k)+(-1)^\beta 2\beta -1
\end{equation}

is non-negative for $q<k<2q-2$.
\end{lemma}

\begin{proof}

We proceed by cases on the parity of $\alpha$ and $\beta$. Throughout, we freely use the identity \[\left\lfloor{x\over 2}\right\rfloor=\left\lceil {x+1\over 2}\right\rceil-1\]

\noindent\textbf{Case 1:}
We start with $\alpha$ being even. In this case, while we originally had the assumption $k<2q-2$, we may  $k\leq 2q-4$, as if $k=2q-3$, we have 
\[\alpha=\left\lfloor {2q-3\over 2}\right\rfloor=\lfloor q-{3\over 2}\rfloor=q-2.\].

\noindent \textbf{Case 1(a):}
Now, if  $\beta$ is additionally even, Equation \ref{eq:floorsnonneg} becomes

\begin{align*}
-\alpha+2q-k+2\beta-1&=\left\lceil{-k\over 2}-k \right\rceil+ 2\left\lceil q+{k-q+1\over 2}\right\rceil-3\\
&=\left\lceil{-3k\over 2} \right\rceil+ 2\left\lceil {k+q+1\over 2}\right\rceil-3.
\end{align*}

Observe $-3k$ is even if and only if $k+q+1$ is even, as $q$ is an odd prime. When both are even, both values in the ceiling functions are integers, and we get ${2q+2-k\over 2}-3$, and since we may assume $k\leq 2q-4$, this becomes
\[{2q+2-k\over 2}-3\geq {6\over 2}-3=0\]
as desired. When both are odd, then $k-q$ is even, and we have
\begin{align*}
\left\lceil{-3k\over 2} \right\rceil+ 2\left\lceil {k+q+1\over 2}\right\rceil-3&=
\left\lceil{-3k\over 2} \right\rceil+ k+q+2\left\lceil {1\over 2}\right\rceil-3\\
&=
\left\lceil{2q-k\over 2} \right\rceil-1,
\end{align*}

Which must be nonnegative as $k<2q-2$. 

\noindent\textbf{Case 1(b):} Suppose $\beta$ is odd. Then Equation \ref{eq:floorsnonneg} becomes
\begin{align*}
k-\alpha -2\beta-1&=\left\lceil {k\over 2}\right\rceil+2\left\lceil{q-k\over 2}\right\rceil-1\\
&= \left\lceil {q\over 2}\right\rceil+\left\lceil{q-k\over 2}\right\rceil-1,
\end{align*}
where the last equality holds as $k$ is even if and only if $q-k$ is odd, and so one of ${k\over 2}$ or $q-k\over 2$ is always an integer. As at the end of Case 1(a), we proceed in two ways. Why $k$ is odd, ${q-k\over 2}$ is an integer and the above becomes
\[\left\lfloor{2q-k\over 2}\right\rfloor -1\]
which is nonnegative provided $k\leq 2q-2$. Otherwise, when $k$ is even, the above becomes
\[2\left\lfloor{q\over 2}\right\rfloor -{k\over 2}-1=q-{k\over 2}\geq q-{2q-2\over 2}=1.\]

\noindent \textbf{Case 2:} Now we assume that $\alpha$ is odd. 

\noindent\textbf{Case 2(a):}
When $\beta$ is even, Equation \ref{eq:floorsnonneg} becomes

\begin{align*}
\alpha+2q-2k+2\beta-1
&=\left\lfloor{-3k\over 2}\right\rfloor +2\left\lfloor {k+q\over 2}\right\rfloor-1\\
&=\left\lfloor{2q-k\over 2}\right\rfloor-1
\end{align*}
where the last equality holds since $-3k$ is even if and only if $k+q$ is odd. Since $k<2q-2$, this is always nonnegative. 

\noindent \textbf{Case 2(b):} If $\beta$ is odd, we have 
\begin{align*}
\alpha-2\beta-1&=\alpha+2\left\lceil {q-k\over 2}\right\rceil -1\\
&=\left\lceil{k+1\over 2}\right\rceil+2\left\lceil {q-k\over 2}\right\rceil -2.
\end{align*}
Note $k+1$ is even if and only if $q-k$ is even. When both are even, the fractions in the above ceiling functions are integers, and we get
\[{2q-k+1\over 2}-2\geq {4\over 2}-2=0\]
as $k\leq 2q-3$. If both $k+1$ and $q-k$ are odd, and thus $k$ is even, we have



\begin{align*}
\left\lceil{k+1\over 2}\right\rceil+2\left\lceil {q-k\over 2}\right\rceil -2&=\left\lceil {1\over 2}\right\rceil+\left\lceil {q-k\over 2}\right \rceil+\left\lceil {q\over 2}\right \rceil-2\\
&=\left\lceil {q-k\over 2}\right \rceil+\left\lceil{q+1\over 2}\right\rceil -1\\
&=\left\lceil {2q-k+1\over 2}\right \rceil -1,
\end{align*}
where here we leverage the fact that $q$ is odd. This is nonnegative since $k<2q-2$.



\end{proof}

\end{document}
